%% Document LaTeX pour le sujet de DS (Devoir Surveillé)
%% Créé par Vibe Code

\documentclass[12pt, a4paper]{article}
\usepackage[utf8]{inputenc}
\usepackage[french]{babel}
\usepackage{amsmath}
\usepackage{amsfonts}
\usepackage{amssymb}
\usepackage{graphicx}
\usepackage{enumitem}
\usepackage{geometry}

%% Configuration de la page
\geometry{
    left=2cm,
    right=2cm,
    top=2cm,
    bottom=2cm
}

%% Titre et informations
\title{Devoir Surveillé}
\author{}
\date{\today}

%% Début du document
\begin{document}

\maketitle

\section*{Consignes}
\begin{itemize}
    \item Durée : 2 heures
    \item Documents autorisés : Aucun
    \item Calculatrice autorisée : Oui (selon les règles de l'établissement)
    \item Répondre directement sur la feuille ou sur une copie séparée.
\end{itemize}

\section*{Exercice 1}
\subsection*{Énoncé}
Résoudre les questions suivantes :

\begin{enumerate}
    \item Calculer la dérivée de la fonction $f(x) = x^2 + 3x - 5$.
    \item Résoudre l'équation $2x^2 - 4x + 2 = 0$.
    \item Simplifier l'expression $\frac{2x + 4}{x^2 - 4}$.
\end{enumerate}

\subsection*{Barème}
\begin{itemize}
    \item Question 1 : 3 points
    \item Question 2 : 3 points
    \item Question 3 : 4 points
\end{itemize}

\section*{Exercice 2}
\subsection*{Énoncé}
Soit la fonction $g(x) = \sin(x) + \cos(x)$. 

\begin{enumerate}
    \item Calculer $g(0)$ et $g(\frac{\pi}{2})$. 
    \item Déterminer la dérivée de $g(x)$. 
    \item Étudier les variations de $g(x)$ sur l'intervalle $[0, \pi]$. 
\end{enumerate}

\subsection*{Barème}
\begin{itemize}
    \item Question 1 : 2 points
    \item Question 2 : 3 points
    \item Question 3 : 5 points
\end{itemize}

\section*{Exercice 3 (Bonus)}
\subsection*{Énoncé}
Démontrer que pour tout $x \in \mathbb{R}$, on a :
\begin{equation*}
    \sin^2(x) + \cos^2(x) = 1
\end{equation*}

\subsection*{Barème}
\begin{itemize}
    \item Question : 5 points (bonus)
\end{itemize}

\end{document}
